\documentclass[10pt,a4paper]{amsart}
\usepackage{amsmath,amssymb,amsthm,graphicx}
\usepackage{verbatim}
\newenvironment{reference}{\comment}{\endcomment}
%\newenvironment{reference}{}{}
\newenvironment{slogan}{\comment}{\endcomment}
\newenvironment{history}{\comment}{\endcomment}

% For commutative diagrams we use Xy-pic
\usepackage[all]{xy}

% We use 2cell for 2-commutative diagrams.
\xyoption{2cell}
% Theorem environments.
%
\theoremstyle{plain}
\newtheorem{theorem}[subsection]{Theorem}
\newtheorem{proposition}[subsection]{Proposition}
\newtheorem{lemma}[subsection]{Lemma}

\theoremstyle{definition}
\newtheorem{definition}[subsection]{Definition}
\newtheorem{example}[subsection]{Example}
\newtheorem{exercise}[subsection]{Exercise}
\newtheorem{situation}[subsection]{Situation}

\theoremstyle{remark}
\newtheorem{remark}[subsection]{Remark}
\newtheorem{remarks}[subsection]{Remarks}

\numberwithin{equation}{subsection}

% Macros
%
\def\lim{\mathop{\mathrm{lim}}\nolimits}
\def\colim{\mathop{\mathrm{colim}}\nolimits}
\def\Spec{\mathop{\mathrm{Spec}}}
\def\Hom{\mathop{\mathrm{Hom}}\nolimits}
\def\Ext{\mathop{\mathrm{Ext}}\nolimits}
\def\SheafHom{\mathop{\mathcal{H}\!\mathit{om}}\nolimits}
\def\SheafExt{\mathop{\mathcal{E}\!\mathit{xt}}\nolimits}
\def\Sch{\mathit{Sch}}
\def\Mor{\mathop{\mathrm{Mor}}\nolimits}
\def\Ob{\mathop{\mathrm{Ob}}\nolimits}
\def\Sh{\mathop{\mathit{Sh}}\nolimits}
\def\NL{\mathop{N\!L}\nolimits}
\def\CH{\mathop{\mathrm{CH}}\nolimits}
\def\proetale{{pro\text{-}\acute{e}tale}}
\def\etale{{\acute{e}tale}}
\def\QCoh{\mathit{QCoh}}
\def\Ker{\mathop{\mathrm{Ker}}}
\def\Im{\mathop{\mathrm{Im}}}
\def\Coker{\mathop{\mathrm{Coker}}}
\def\Coim{\mathop{\mathrm{Coim}}}

% Boxtimes
%
\DeclareMathSymbol{\boxtimes}{\mathbin}{AMSa}{"02}

%
% Macros for moduli stacks/spaces
%
\def\QCohstack{\mathcal{QC}\!\mathit{oh}}
\def\Cohstack{\mathcal{C}\!\mathit{oh}}
\def\Spacesstack{\mathcal{S}\!\mathit{paces}}
\def\Quotfunctor{\mathrm{Quot}}
\def\Hilbfunctor{\mathrm{Hilb}}
\def\Curvesstack{\mathcal{C}\!\mathit{urves}}
\def\Polarizedstack{\mathcal{P}\!\mathit{olarized}}
\def\Complexesstack{\mathcal{C}\!\mathit{omplexes}}
% \Pic is the operator that assigns to X its picard group, usage \Pic(X)
% \Picardstack_{X/B} denotes the Picard stack of X over B
% \Picardfunctor_{X/B} denotes the Picard functor of X over B
\def\Pic{\mathop{\mathrm{Pic}}\nolimits}
\def\Picardstack{\mathcal{P}\!\mathit{ic}}
\def\Picardfunctor{\mathrm{Pic}}
\def\Deformationcategory{\mathcal{D}\!\mathit{ef}}
\usepackage{cleveref}
\crefname{originallemma}{original lemma}{original lemmas}
\Crefname{originallemma}{Original lemma}{Original lemmas}
\makeatletter
\newlabel{duality-lemma-gorenstein}{{0BFQ}{source}{}{sourcepointer-1}{}}
\newlabel{duality-lemma-gorenstein@cref}{{[originallemma][0][]0BFQ}{[1][0][]source}{}{}{}}
\makeatother
\crefname{originalreference}{original reference}{original references}
\Crefname{originalreference}{Original reference}{Original references}
\makeatletter
\newlabel{lemma-closed-immersion-dim-1-CM}{{0E33}{source}{}{sourcepointer-2}{}}
\newlabel{lemma-closed-immersion-dim-1-CM@cref}{{[originalreference][0][]0E33}{[1][0][]source}{}{}{}}
\makeatother
\crefname{originalreference}{original reference}{original references}
\Crefname{originalreference}{Original reference}{Original references}
\makeatletter
\newlabel{lemma-duality-dim-1}{{0BS2}{source}{}{sourcepointer-3}{}}
\newlabel{lemma-duality-dim-1@cref}{{[originalreference][0][]0BS2}{[1][0][]source}{}{}{}}
\makeatother
\crefname{originallemma}{original lemma}{original lemmas}
\Crefname{originallemma}{Original lemma}{Original lemmas}
\makeatletter
\newlabel{morphisms-lemma-scheme-theoretic-union}{{0C4J}{source}{}{sourcepointer-4}{}}
\newlabel{morphisms-lemma-scheme-theoretic-union@cref}{{[originallemma][0][]0C4J}{[1][0][]source}{}{}{}}
\makeatother
\title{Duality exact sequences for a reduced closed subscheme of a Gorenstein curve and its scheme-theoretic complement}
\author{The Stacks Project Authors}
\date{stacks-project:curves.tex:local-snapshot-20261012; GNU FDL 1.2 or later}
\begin{document}
\maketitle
\noindent\textit{Editorial note.} All delivered mathematical text is the author's own, copied from the pinned Stacks Project chapter file curves.tex; no statement or proof step was rewritten, reordered or removed.\par\medskip
\noindent\textit{Editorial note.} The counted proof is the author's complete argument: the union sequence, the four kernel descriptions with the diagram chase, the six structure-sheaf exact sequences, the reduction by symmetry, and the closing annihilator identification by evaluation at 1 are all retained verbatim.\par\medskip
\noindent\textit{Editorial note.} The label lemma-closed-subscheme-reduced-gorenstein names the counted statement, and the proof environment immediately following it in the source is delivered verbatim as its proof.\par\medskip
\noindent\textit{Editorial note.} The counted statement is carried by the author's own lemma environment in the chapter source.\par\medskip
\noindent\textit{Allowed standard prerequisites (editorial).}\begin{itemize}
\item Scheme-theoretic unions and the union sequence. Morphisms, Lemma morphisms-lemma-scheme-theoretic-union. Supplies 0 -> O\_X -> O\_Y (+) O\_Z -> O\_\{Y cap Z\} -> 0 for X = Y  union  Z and the kernel descriptions of the ideal sheaves used in the diagram chase.
\item Duality for 1-dimensional CM closed subschemes. Same chapter, Lemma lemma-duality-dim-1 for the dualizing modules and Lemma lemma-closed-immersion-dim-1-CM; Duality for Schemes, Lemma duality-lemma-gorenstein. Makes omega\_X invertible on the Gorenstein X, identifies i\_*omega\_Y as SheafHom(i\_*O\_Y, omega\_X), and trivializes omega\_X over the 0-dimensional intersection.
\end{itemize}\medskip
\ifdefined\hypertarget\hypertarget{sourcepointer-1}{}\fi
\noindent\textit{Source reference.} Original lemma 0BFQ: sources/stacks-chow/tags.txt, source lines 14393--14393. Prerequisite pointer invoked at proof line(s) 1081; it names a standard fact of a sibling chapter that the counted argument uses, and every step of the counted proof built on it is retained verbatim in the delivered body.\par\medskip
\ifdefined\hypertarget\hypertarget{sourcepointer-2}{}\fi
\noindent\textit{Source reference.} Original reference 0E33: sources/stacks-curves/pinned-original-source.tex, source lines 976--976. Prerequisite pointer invoked at proof line(s) 1093; it names a standard fact of the same chapter that the counted argument uses, and every step of the counted proof built on it is retained verbatim in the delivered body.\par\medskip
\ifdefined\hypertarget\hypertarget{sourcepointer-3}{}\fi
\noindent\textit{Source reference.} Original reference 0BS2: sources/stacks-curves/pinned-original-source.tex, source lines 796--796. Prerequisite pointer invoked at proof line(s) 1048; it names a standard fact of the same chapter that the counted argument uses, and every step of the counted proof built on it is retained verbatim in the delivered body.\par\medskip
\ifdefined\hypertarget\hypertarget{sourcepointer-4}{}\fi
\noindent\textit{Source reference.} Original lemma 0C4J: sources/stacks-chow/tags.txt, source lines 15229--15229. Prerequisite pointer invoked at proof line(s) 1055; it names a standard fact of a sibling chapter that the counted argument uses, and every step of the counted proof built on it is retained verbatim in the delivered body.\par\medskip
\begin{lemma}
\label{lemma-closed-subscheme-reduced-gorenstein}
Let $X$ be a proper scheme over a field $k$ which is
Gorenstein, reduced, and equidimensional of dimension $1$.
Let $i : Y \to X$ be a reduced closed subscheme equidimensional
of dimension $1$. Let $j : Z \to X$ be the scheme theoretic
closure of $X \setminus Y$. Then
\begin{enumerate}
\item $Y$ and $Z$ are Cohen-Macaulay,
\item if $\mathcal{I} \subset \mathcal{O}_X$,
resp.\ $\mathcal{J} \subset \mathcal{O}_X$ is the ideal sheaf of
$Y$, resp.\ $Z$ in $X$, then
$$
\mathcal{I} = i_*\mathcal{I}'
\quad\text{and}\quad
\mathcal{J} = j_*\mathcal{J}'
$$
where $\mathcal{I}' \subset \mathcal{O}_Z$,
resp.\ $\mathcal{J}' \subset \mathcal{O}_Y$ is the ideal sheaf
of $Y \cap Z$ in $Z$, resp.\ $Y$,
\item $\omega_Y = \mathcal{J}'(i^*\omega_X)$ and
$i_*(\omega_Y) = \mathcal{J}\omega_X$,
\item $\omega_Z = \mathcal{I}'(i^*\omega_X)$ and
$i_*(\omega_Z) = \mathcal{I}\omega_X$,
\item we have the following short exact sequences
\begin{align*}
0 \to \omega_X \to i_*i^*\omega_X \oplus j_*j^*\omega_X \to
\mathcal{O}_{Y \cap Z} \to 0 \\
0 \to i_*\omega_Y \to \omega_X \to j_*j^*\omega_X \to 0 \\
0 \to j_*\omega_Z \to \omega_X \to i_*i^*\omega_X \to 0 \\
0 \to i_*\omega_Y \oplus j_*\omega_Z \to \omega_X \to
\mathcal{O}_{Y \cap Z} \to 0 \\
0 \to \omega_Y \to i^*\omega_X \to \mathcal{O}_{Y \cap Z} \to 0 \\
0 \to \omega_Z \to j^*\omega_X \to \mathcal{O}_{Y \cap Z} \to 0
\end{align*}
\end{enumerate}
Here $\omega_X$, $\omega_Y$, $\omega_Z$ are as in
Lemma \ref{lemma-duality-dim-1}.
\end{lemma}

\begin{proof}
A reduced $1$-dimensional Noetherian scheme is Cohen-Macaulay, so
(1) is true. Since $X$ is reduced, we see that $X = Y \cup Z$
scheme theoretically. With notation as in
Morphisms, Lemma \ref{morphisms-lemma-scheme-theoretic-union}
and by the statement of that lemma
we have a short exact sequence
$$
0 \to \mathcal{O}_X \to
\mathcal{O}_Y \oplus \mathcal{O}_Z \to \mathcal{O}_{Y \cap Z} \to 0
$$
Since $\mathcal{J} = \Ker(\mathcal{O}_X \to \mathcal{O}_Z)$,
$\mathcal{J}' = \Ker(\mathcal{O}_Y \to \mathcal{O}_{Y \cap Z})$,
$\mathcal{I} = \Ker(\mathcal{O}_X \to \mathcal{O}_Y)$, and
$\mathcal{I}' = \Ker(\mathcal{O}_Z \to \mathcal{O}_{Y \cap Z})$
a diagram chase implies (2).
Observe that $\mathcal{I} + \mathcal{J}$ is the ideal sheaf
of $Y \cap Z$ and that $\mathcal{I} \cap \mathcal{J} = 0$.
Hence we have the following exact sequences
\begin{align*}
0 \to \mathcal{O}_X \to \mathcal{O}_Y \oplus \mathcal{O}_Z \to
\mathcal{O}_{Y \cap Z} \to 0 \\
0 \to \mathcal{J} \to \mathcal{O}_X \to \mathcal{O}_Z \to 0 \\
0 \to \mathcal{I} \to \mathcal{O}_X \to \mathcal{O}_Y \to 0 \\
0 \to \mathcal{J} \oplus \mathcal{I} \to \mathcal{O}_X \to
\mathcal{O}_{Y \cap Z} \to 0 \\
0 \to \mathcal{J}' \to \mathcal{O}_Y \to \mathcal{O}_{Y \cap Z} \to 0 \\
0 \to \mathcal{I}' \to \mathcal{O}_Z \to \mathcal{O}_{Y \cap Z} \to 0
\end{align*}
Since $X$ is Gorenstein $\omega_X$ is an invertible $\mathcal{O}_X$-module
(Duality for Schemes, Lemma \ref{duality-lemma-gorenstein}).
Since $Y \cap Z$ has dimension $0$ we have
$\omega_X|_{Y \cap Z} \cong \mathcal{O}_{Y \cap Z}$.
Thus if we prove (3) and (4), then we obtain the short exact
sequences of the lemma by tensoring the above
short exact sequence with the invertible module $\omega_X$.
By symmetry it suffices to prove (3) and by
(2) it suffices to prove $i_*(\omega_Y) = \mathcal{J}\omega_X$.

\medskip\noindent
We have
$i_*\omega_Y = \SheafHom_{\mathcal{O}_X}(i_*\mathcal{O}_Y, \omega_X)$
by Lemma \ref{lemma-closed-immersion-dim-1-CM}.
Again using that $\omega_X$ is invertible
we finally conclude that it suffices to show
$\SheafHom_{\mathcal{O}_X}(\mathcal{O}_X/\mathcal{I}, \mathcal{O}_X)$
maps isomorphically to $\mathcal{J}$ by evaluation at $1$.
In other words, that $\mathcal{J}$ is the annihilator of
$\mathcal{I}$. This follows from the above.
\end{proof}

\end{document}
